\begin{thema}{Δ}
Ο πληθυσμός βακτηρίων (σε χιλιάδες) τη χρονική στιγμή $t \ge 0$ (σε ώρες) περιγράφεται από τη συνάρτηση:
\[ P(t) = \begin{cases} 1000 e^{t/2} & 0 \le t \le 10 \\ A e^{-t/2} & t > 10 \end{cases} \]
όπου $A > 0$ σταθερά.
\begin{erwthma}
    \item Να βρείτε την τιμή της σταθεράς $A$ ώστε η $P$ να είναι συνεχής στο $t = 10$.
    \item Να εξετάσετε αν η $P$ είναι παραγωγίσιμη στο $t = 10$ και να ερμηνεύσετε γεωμετρικά το αποτέλεσμα.
    \item Να βρείτε τη χρονική στιγμή $t_0 > 10$ κατά την οποία ο πληθυσμός επιστρέφει στα 1000 (χιλιάδες βακτήρια). Εκφράστε το αποτέλεσμα ως $t_0 = 10 + k$ και βρείτε το $k$.
    \item Να υπολογίσετε το ολοκλήρωμα $\displaystyle\dint_0^{10} P(t)\,dt$, το οποίο εκφράζει τη «συνολική έκθεση» στα βακτήρια κατά τη φάση ανάπτυξης. Χωρίς υπολογισμό, να συγκρίνετε αυτό με το $\displaystyle\dint_{10}^{t_0} P(t)\,dt$ αξιοποιώντας μια κατάλληλη σχέση συμμετρίας για τις τιμές του $P$.
\end{erwthma}
\end{thema}
